\section*{Apêndice A -- Estudo bibliométrico: informações gerais, rede de coocorrência e mapa temático}
\addcontentsline{toc}{section}{Apêndice A -- Estudo bibliométrico}

\begin{table}[htbp]
\centering
\caption{Estudo bibliométrico: informações principais do corpus (\textit{bibliometrix})}
\label{tab:biblio}
\scriptsize
\begin{tabular}{@{}p{0.47\textwidth}@{\hspace{0.04\textwidth}}p{0.47\textwidth}@{}}
\begin{tabular}[t]{lr}
\toprule
\multicolumn{2}{l}{\textit{(a) Informações gerais}}\\
\midrule
Período coberto & 1952--2025\\
Documentos (artigos em periódicos) & 1992\\
Periódicos & 1109\\
Autores & 3107\\
Taxa de crescimento anual (\%) & 7,26\\
Citações médias por documento & 18,33\\
Coautores por documento & 1,94\\
Coautoria internacional (\%) & 8,333\\
\midrule
\multicolumn{2}{l}{\textit{(b) Periódicos mais produtivos}}\\
\midrule
Public Choice & 54\\
Public Finance Review & 28\\
International Tax and Public Finance & 26\\
Publius: the Journal of Federalism & 24\\
Public Budgeting \& Finance & 23\\
\bottomrule
\end{tabular}
&
\begin{tabular}[t]{lr}
\toprule
\multicolumn{2}{l}{\textit{(c) Países (afiliação dos autores, ocorrências)}}\\
\midrule
Indonésia & 551\\
Estados Unidos & 532\\
Brasil & 461\\
Ucrânia & 126\\
China & 110\\
Reino Unido & 85\\
\midrule
\multicolumn{2}{l}{\textit{(d) Palavras-chave mais frequentes}}\\
\midrule
intergovernmental transfers & 1094\\
fiscal federalism & 541\\
fiscal decentralization & 539\\
flypaper effect & 473\\
local government finance & 375\\
fiscal transfers & 368\\
fiscal illusion & 295\\
local government spending & 255\\
\bottomrule
\end{tabular}
\end{tabular}
\par\smallskip
\begin{minipage}{0.97\textwidth}\scriptsize
Fonte: elaboração própria com o pacote \textit{bibliometrix} \cite{aria2017} aplicado a registros do OpenAlex (consultas em título e resumo em inglês, português e espanhol, descritas na nota da Tabela~\ref{tab:baker}; apenas artigos com autoria e fonte primária classificada como periódico). Código: \texttt{05\_bibliometria.R}.
\end{minipage}
\end{table}


\begin{figure}[htbp]
\centering
\caption{Rede de coocorrência das 40 palavras-chave mais frequentes}\label{fig:rede}
\includegraphics[width=0.85\textwidth]{rede_coocorrencia.png}
\par\smallskip
\begin{minipage}{0.85\textwidth}\scriptsize
Fonte: elaboração própria com \textit{bibliometrix} (\texttt{biblioNetwork} e \texttt{networkPlot}, algoritmo de Fruchterman-Reingold), corpus do OpenAlex.
\end{minipage}
\end{figure}

\begin{figure}[htbp]
\centering
\caption{Mapa temático (centralidade $\times$ densidade) das palavras-chave}\label{fig:mapa}
\includegraphics[width=0.7\textwidth]{mapa_tematico.png}
\par\smallskip
\begin{minipage}{0.7\textwidth}\scriptsize
Fonte: elaboração própria com \textit{bibliometrix} (\texttt{thematicMap}, 250 termos, frequência mínima de 5).
\end{minipage}
\end{figure}

\section*{Apêndice B -- Repositório, reprodução e fontes}
\addcontentsline{toc}{section}{Apêndice B -- Repositório e reprodução}

O repositório que acompanha este artigo está organizado de acordo com os itens exigidos: \texttt{01\_base\_original} (base do trabalho anterior, sem alterações); \texttt{02\_dados\_brutos\_atualizados} (respostas das APIs do Siconfi, do IBGE, do Eurostat e do OpenAlex em JSON e as mesmas tabelas em CSV e XLSX editáveis); \texttt{03\_dados\_tratados} (bases usadas nas análises, em CSV e XLSX); \texttt{04\_dicionario\_variaveis} (conceito, unidade, fonte, endereço, período, cobertura, tratamento de ausentes e fórmula de cada variável); \texttt{05\_registro\_alteracoes} (comparação entre a base anterior e a nova versão); \texttt{06\_codigos\_e\_planilhas} (scripts Python e R e planilha com fórmulas que reproduzem tratamentos, cálculos, tabelas, gráficos e resultados); \texttt{07\_resultados} (tabelas, figuras e saídas do \textit{bibliometrix}); e \texttt{08\_artigo} (código-fonte deste artigo, em que todos os números do texto são importados automaticamente dos resultados). As instruções de reprodução estão no arquivo \texttt{README.md}.

Os endereços das fontes utilizadas são: Siconfi/DCA, \url{https://apidatalake.tesouro.gov.br/ords/siconfi/tt/dca}; IBGE/Sidra, \url{https://servicodados.ibge.gov.br/api/v3/agregados}; Eurostat, \url{https://ec.europa.eu/eurostat/api/dissemination/statistics/1.0/data/}; OpenAlex, \url{https://api.openalex.org}.
